\subsubsection{Cascade Bounds and Restricted Influence Structures}
\label{cascade-bounds-prior-art}

The stochastic-reference direction introduces additional close prior art.
\citet{abbe2017} derive cascade bounds using FKG inequalities and
nonbacktracking message passing. The expanded journal treatment
\citep{abbe2020} defines the activation probability of each node before
summing those probabilities into seed-set influence. Its upper construction
combines length-indexed incoming messages by complement product
(Section~3.1, Equations~6--8). Its lower construction processes a directed
acyclic subnetwork in topological order with an incoming-message inequality
(Section~4.1, Equations~21--23), including an $O(|V|+|E|)$ implementation.
Thus nodewise bounds, positive-association arguments, and efficient DAG
lower bounds already have clear precedents; the aggregate influence target
does not erase their relevance to receiving-node evaluation.

The present benchmark uses independently sampled local activations, a
single topological complement-product pass for $R_j$, a recursive maximum
lower bound, and an optional target-specific restricted tree. Abbe et al.'s
upper messages are indexed by path length, while their lower recurrence
and subnetwork restriction differ from these two proposed lower
comparators. An auxiliary certain seed with independent edges of weights
$B_j$ can encode the random local starts, so that input convention alone
will not be presented as a novelty distinction. The published DAG lower
recursion will be tested as a separate matched nodewise comparator after
this transformation in Section~\ref{reference-comparators}. No ordering
between its lower values and $M_j,T_j$ will be assumed. The intended result
is a transparent DAG specialization and
measured tightness, ranking, and computational tradeoffs under this
declared benchmark; any stronger novelty claim will depend on the final
comparison and evidence.

\citet{chen2010mia} construct maximum-influence in-arborescences by retaining
consistent maximum-probability paths and applying an influence threshold.
Their Algorithm 2 evaluates activation through complement-product recursion
on a restricted tree, where multiple seed influences can be aggregated.
The tree removes the dependency issue associated with alternative paths
from the same origin. MIA therefore provides both a mathematical precedent
and a relevant restriction-based comparator; it is not equivalent to
retaining one largest contribution globally. Its complete algorithm selects
seeds for influence maximization, whereas this study fixes local inputs and
evaluates receiving-state queries.

\citet{chen2010ldag} instead develop local DAG influence computation under
the Linear Threshold model. Their DAG activation recurrence is additive
under that model's incoming-weight constraints. LDAG is consequently
related structural approximation literature, not an interchangeable
Independent Cascade reference or the same ancestry-grouping algorithm.

These studies will be distinguished from the prospective diagnostic that
groups overlapping local-origin label sets and keeps one group maximum.
Overlap grouping does not construct target-specific arborescences and does
not automatically supply exact restricted-process probabilities. The study
will compare its assumptions and failure cases with the tree restriction
before any extension is described as original. A targeted literature review
of equivalent approximations will continue; neither the bracket proof nor
the proposed evaluation protocol is a verified priority claim.
